% =====================================================================
% Paper 11 — Λ cancellation PROGRAMME note (NOT a Pillar-11 closure paper)
% Status: [DRAFT] [NEEDS_UPDATE] [INTERNAL_INCONSISTENCY_REPAIRED]
%   2026-05-02: paper internally inconsistent --- abstract claimed
%   "verification within F-GAP gate |Lambda|/M_P^4 < 10^{-120}", but
%   §Numerical-result body computes Lambda/M_P^4 ~ 10^{-118}..10^{-119},
%   which DOES NOT PASS the stated gate. The paper failed its own
%   pre-registered falsification gate per CLAUDE.md §6.3.4 quantitative
%   sanity check. Abstract and Numerical-result section now flag this
%   internal inconsistency explicitly; Pillar-11 closure claim retracted.
%   Recast as Λ cancellation PROGRAMME note pending verdict-shell
%   resolution.
% AUDIT-FLAG: AUDIT-2026-05-02-Wave7-Aux-Epoch-Overclaim
%   (Math314-AddD extension covering Wave 1/4/5 papers Paper-09..16).
% Compile: pdflatex Paper-11.tex (×2 for refs)
% Canonical source: Math58-v7, Math58-v7-Addendum-A, Math58-v7-AddC,
%                   Math314-AddD
% =====================================================================
% --- TECT canonical preamble (see _shared/tect-preamble.tex) ----------
% =====================================================================
% TECT canonical LaTeX preamble (revtex4-2 / single-column / theorem-safe)
% =====================================================================
% Maintainer: Jusang Lee (jtkor@outlook.com)
% Binding from: 2026-05-06
% Purpose: a single source-of-truth preamble for every TECT paper
% (Paper-00 .. Paper-10 + Auxiliary notes). Designed to (a) eliminate
% font-corruption on pdflatex (T1 fontenc + lmodern), (b) make
% theorem-heavy mathematical-physics content render cleanly in a
% single-column layout (PRD class option, NOT PRL two-column), (c)
% provide consistent theorem numbering across all TECT papers via
% amsthm with a shared counter set, (d) make hyperref + cleveref
% cross-references resolve cleanly with the standard two-pass
% pdflatex compile.
%
% Usage in each Paper-NN.tex:
%   % =====================================================================
% TECT canonical LaTeX preamble (revtex4-2 / single-column / theorem-safe)
% =====================================================================
% Maintainer: Jusang Lee (jtkor@outlook.com)
% Binding from: 2026-05-06
% Purpose: a single source-of-truth preamble for every TECT paper
% (Paper-00 .. Paper-10 + Auxiliary notes). Designed to (a) eliminate
% font-corruption on pdflatex (T1 fontenc + lmodern), (b) make
% theorem-heavy mathematical-physics content render cleanly in a
% single-column layout (PRD class option, NOT PRL two-column), (c)
% provide consistent theorem numbering across all TECT papers via
% amsthm with a shared counter set, (d) make hyperref + cleveref
% cross-references resolve cleanly with the standard two-pass
% pdflatex compile.
%
% Usage in each Paper-NN.tex:
%   % =====================================================================
% TECT canonical LaTeX preamble (revtex4-2 / single-column / theorem-safe)
% =====================================================================
% Maintainer: Jusang Lee (jtkor@outlook.com)
% Binding from: 2026-05-06
% Purpose: a single source-of-truth preamble for every TECT paper
% (Paper-00 .. Paper-10 + Auxiliary notes). Designed to (a) eliminate
% font-corruption on pdflatex (T1 fontenc + lmodern), (b) make
% theorem-heavy mathematical-physics content render cleanly in a
% single-column layout (PRD class option, NOT PRL two-column), (c)
% provide consistent theorem numbering across all TECT papers via
% amsthm with a shared counter set, (d) make hyperref + cleveref
% cross-references resolve cleanly with the standard two-pass
% pdflatex compile.
%
% Usage in each Paper-NN.tex:
%   \input{../_shared/tect-preamble.tex}
%   \begin{document}
%   ...
%   \end{document}
%
% Build (every paper): `pdflatex Paper-NN && pdflatex Paper-NN`
% (the second pass resolves \ref{} / \cref{} forward references).
% A bash/PowerShell wrapper that automates this is at
% Docs/papers/papers/_shared/build-paper.sh / build-paper.ps1.
% =====================================================================

% ---- Document class --------------------------------------------------
% PRD class option of revtex4-2 with [reprint,onecolumn] gives the same
% APS heading / by-line / abstract style as PRL but in a wider single
% column that handles long display equations and theorem environments
% without column-break artefacts. nofootinbib keeps citations out of
% footnotes. The 11pt option restores standard font metrics; with T1
% fontenc + lmodern this gives non-bitmap glyphs.
\documentclass[%
  aps,prd,reprint,onecolumn,nofootinbib,11pt,%
  amsmath,amssymb,floatfix,longbibliography%
]{revtex4-2}

% ---- Encoding & fonts ------------------------------------------------
\usepackage[T1]{fontenc}        % proper 8-bit font encoding (no font corruption)
\usepackage[utf8]{inputenc}     % allow UTF-8 source
\usepackage{lmodern}            % scalable Latin Modern (replaces bitmap CM)
\usepackage{textcomp}           % extra text symbols
\usepackage{microtype}          % subtle kerning improvements

% ---- UTF-8 → LaTeX character mappings --------------------------------
% Maps the Unicode characters that occur in TECT body text (legacy from
% the chat-content auto-archival rule, CLAUDE.md §4) to LaTeX-safe
% commands. Without these declarations, T1+inputenc rejects the chars
% with "Unicode character X not set up for use with LaTeX".
% Adding declarations centrally (here) is preferable to per-file edits.
\DeclareUnicodeCharacter{00A7}{\S}              % §  (section sign)
\DeclareUnicodeCharacter{00B2}{\ensuremath{^{2}}}        % ²
\DeclareUnicodeCharacter{00B3}{\ensuremath{^{3}}}        % ³
\DeclareUnicodeCharacter{00B5}{\ensuremath{\mu}}         % µ (micro)
\DeclareUnicodeCharacter{00B7}{\ensuremath{\cdot}}       % · (middle dot)
\DeclareUnicodeCharacter{00D7}{\ensuremath{\times}}      % ×
\DeclareUnicodeCharacter{00E4}{\"a}             % ä
\DeclareUnicodeCharacter{00E9}{\'e}             % é
\DeclareUnicodeCharacter{00F6}{\"o}             % ö
\DeclareUnicodeCharacter{00FC}{\"u}             % ü
\DeclareUnicodeCharacter{010C}{\v{C}}           % Č
\DeclareUnicodeCharacter{010D}{\v{c}}           % č
\DeclareUnicodeCharacter{0394}{\ensuremath{\Delta}}      % Δ
\DeclareUnicodeCharacter{03A3}{\ensuremath{\Sigma}}      % Σ
\DeclareUnicodeCharacter{03A9}{\ensuremath{\Omega}}      % Ω
\DeclareUnicodeCharacter{03B1}{\ensuremath{\alpha}}      % α
\DeclareUnicodeCharacter{03B2}{\ensuremath{\beta}}       % β
\DeclareUnicodeCharacter{03B3}{\ensuremath{\gamma}}      % γ
\DeclareUnicodeCharacter{03B4}{\ensuremath{\delta}}      % δ
\DeclareUnicodeCharacter{03B5}{\ensuremath{\varepsilon}} % ε
\DeclareUnicodeCharacter{03BA}{\ensuremath{\kappa}}      % κ
\DeclareUnicodeCharacter{03BB}{\ensuremath{\lambda}}     % λ
\DeclareUnicodeCharacter{03BC}{\ensuremath{\mu}}         % μ
\DeclareUnicodeCharacter{03BD}{\ensuremath{\nu}}         % ν
\DeclareUnicodeCharacter{03C0}{\ensuremath{\pi}}         % π
\DeclareUnicodeCharacter{03C1}{\ensuremath{\rho}}        % ρ
\DeclareUnicodeCharacter{03C3}{\ensuremath{\sigma}}      % σ
\DeclareUnicodeCharacter{03C4}{\ensuremath{\tau}}        % τ
\DeclareUnicodeCharacter{03C6}{\ensuremath{\varphi}}     % φ
\DeclareUnicodeCharacter{03C7}{\ensuremath{\chi}}        % χ
\DeclareUnicodeCharacter{03C8}{\ensuremath{\psi}}        % ψ
\DeclareUnicodeCharacter{03C9}{\ensuremath{\omega}}      % ω
\DeclareUnicodeCharacter{2013}{--}              % – (en dash)
\DeclareUnicodeCharacter{2014}{---}             % — (em dash)
\DeclareUnicodeCharacter{2018}{`}               % ‘
\DeclareUnicodeCharacter{2019}{'}               % ’
\DeclareUnicodeCharacter{201C}{``}              % “
\DeclareUnicodeCharacter{201D}{''}              % ”
\DeclareUnicodeCharacter{2070}{\ensuremath{^{0}}}        % ⁰
\DeclareUnicodeCharacter{2071}{\ensuremath{^{i}}}        % ⁱ
\DeclareUnicodeCharacter{2122}{\textsuperscript{TM}}     % ™
\DeclareUnicodeCharacter{2126}{\ensuremath{\Omega}}      % Ω (ohm sign)
\DeclareUnicodeCharacter{210F}{\ensuremath{\hbar}}       % ℏ
\DeclareUnicodeCharacter{2115}{\ensuremath{\mathbb{N}}}  % ℕ
\DeclareUnicodeCharacter{2102}{\ensuremath{\mathbb{C}}}  % ℂ
\DeclareUnicodeCharacter{211D}{\ensuremath{\mathbb{R}}}  % ℝ
\DeclareUnicodeCharacter{2124}{\ensuremath{\mathbb{Z}}}  % ℤ
\DeclareUnicodeCharacter{211A}{\ensuremath{\mathbb{Q}}}  % ℚ
\DeclareUnicodeCharacter{2192}{\ensuremath{\to}}         % →
\DeclareUnicodeCharacter{21D2}{\ensuremath{\Rightarrow}} % ⇒
\DeclareUnicodeCharacter{2200}{\ensuremath{\forall}}     % ∀
\DeclareUnicodeCharacter{2203}{\ensuremath{\exists}}     % ∃
\DeclareUnicodeCharacter{2208}{\ensuremath{\in}}         % ∈
\DeclareUnicodeCharacter{2209}{\ensuremath{\notin}}      % ∉
\DeclareUnicodeCharacter{2227}{\ensuremath{\wedge}}      % ∧
\DeclareUnicodeCharacter{2228}{\ensuremath{\vee}}        % ∨
\DeclareUnicodeCharacter{2229}{\ensuremath{\cap}}        % ∩
\DeclareUnicodeCharacter{222A}{\ensuremath{\cup}}        % ∪
\DeclareUnicodeCharacter{2245}{\ensuremath{\cong}}       % ≅
\DeclareUnicodeCharacter{2248}{\ensuremath{\approx}}     % ≈
\DeclareUnicodeCharacter{2260}{\ensuremath{\neq}}        % ≠
\DeclareUnicodeCharacter{2264}{\ensuremath{\leq}}        % ≤
\DeclareUnicodeCharacter{2265}{\ensuremath{\geq}}        % ≥
\DeclareUnicodeCharacter{22C5}{\ensuremath{\cdot}}       % ⋅
\DeclareUnicodeCharacter{2295}{\ensuremath{\oplus}}      % ⊕
\DeclareUnicodeCharacter{2297}{\ensuremath{\otimes}}     % ⊗

% ---- Mathematics & symbols ------------------------------------------
\usepackage{amsmath,amssymb,bm,mathrsfs,mathtools}
\usepackage{amsthm}             % theorem environments (load BEFORE hyperref)

% ---- Theorem environments (shared counter set, TECT canonical) -------
% All theorems / lemmas / propositions / corollaries / remarks share a
% single counter so cross-references read consistently across a paper.
% Definitions get a separate counter (definitions are numbered in their
% own sequence by editorial convention).
\newtheorem{theorem}{Theorem}
\newtheorem{lemma}[theorem]{Lemma}
\newtheorem{proposition}[theorem]{Proposition}
\newtheorem{corollary}[theorem]{Corollary}

\theoremstyle{remark}
\newtheorem{remark}[theorem]{Remark}
\newtheorem{example}[theorem]{Example}

\theoremstyle{definition}
\newtheorem{definition}[theorem]{Definition}
\newtheorem{conjecture}[theorem]{Conjecture}
\newtheorem{hypothesis}[theorem]{Hypothesis}

% ---- Tables / floats / graphics --------------------------------------
\usepackage{booktabs}           % \toprule / \midrule / \bottomrule
\usepackage{array}
\usepackage{graphicx}
\usepackage{enumitem}           % \begin{itemize}[leftmargin=...] options
\usepackage{hyphenat}           % \hyp{} for breakable hyphens in \texttt
\usepackage{seqsplit}           % \seqsplit{} for long unbreakable strings
\sloppy                         % allow stretchier inter-word spacing to
                                % reduce overfull \hbox warnings on long URLs
                                % and long \texttt tokens (paragraph-level).
\emergencystretch=3em           % last-resort hbox stretch budget

% ---- Cross-references (hyperref last among the standard packages,
%      cleveref AFTER hyperref) ----------------------------------------
\usepackage[%
  colorlinks=true,%
  linkcolor=blue,%
  citecolor=blue,%
  urlcolor=blue,%
  breaklinks=true,%
  bookmarks=true,%
  bookmarksnumbered=true%
]{hyperref}
\usepackage[capitalise,nameinlink]{cleveref}

% ---- TECT-canonical author block macros -----------------------------
% (defined here so every paper has the same author / affiliation style)
\newcommand{\tectauthor}{%
  \author{Jusang Lee}%
  \email{jtkor@outlook.com}%
  \affiliation{Independent researcher, Republic of Korea}%
}

% ---- TECT-canonical archive footer ----------------------------------
\newcommand{\tectarchivefooter}{%
  \par\medskip
  \noindent\small\textit{Canonical archive}:
  \url{https://github.com/TECT-OpenPhysics/TECT}.\par%
}

% ---- TECT TODO / AUDIT-FLAG / HONEST-SCOPE inline marks --------------
% Used in drafts to highlight pending audit items without disturbing
% the body text style. Suppress via \tectsuppressmarks (define empty).
\providecommand{\tectauditflag}[1]{\textcolor{red}{\textbf{[AUDIT-FLAG: #1]}}}
\providecommand{\tecttodo}[1]{\textcolor{orange}{\textbf{[TODO: #1]}}}
\providecommand{\tecthonestscope}[1]{\textit{Honest scope: #1.}}

% =====================================================================
% End of TECT canonical preamble
% =====================================================================

%   \begin{document}
%   ...
%   \end{document}
%
% Build (every paper): `pdflatex Paper-NN && pdflatex Paper-NN`
% (the second pass resolves \ref{} / \cref{} forward references).
% A bash/PowerShell wrapper that automates this is at
% Docs/papers/papers/_shared/build-paper.sh / build-paper.ps1.
% =====================================================================

% ---- Document class --------------------------------------------------
% PRD class option of revtex4-2 with [reprint,onecolumn] gives the same
% APS heading / by-line / abstract style as PRL but in a wider single
% column that handles long display equations and theorem environments
% without column-break artefacts. nofootinbib keeps citations out of
% footnotes. The 11pt option restores standard font metrics; with T1
% fontenc + lmodern this gives non-bitmap glyphs.
\documentclass[%
  aps,prd,reprint,onecolumn,nofootinbib,11pt,%
  amsmath,amssymb,floatfix,longbibliography%
]{revtex4-2}

% ---- Encoding & fonts ------------------------------------------------
\usepackage[T1]{fontenc}        % proper 8-bit font encoding (no font corruption)
\usepackage[utf8]{inputenc}     % allow UTF-8 source
\usepackage{lmodern}            % scalable Latin Modern (replaces bitmap CM)
\usepackage{textcomp}           % extra text symbols
\usepackage{microtype}          % subtle kerning improvements

% ---- UTF-8 → LaTeX character mappings --------------------------------
% Maps the Unicode characters that occur in TECT body text (legacy from
% the chat-content auto-archival rule, CLAUDE.md §4) to LaTeX-safe
% commands. Without these declarations, T1+inputenc rejects the chars
% with "Unicode character X not set up for use with LaTeX".
% Adding declarations centrally (here) is preferable to per-file edits.
\DeclareUnicodeCharacter{00A7}{\S}              % §  (section sign)
\DeclareUnicodeCharacter{00B2}{\ensuremath{^{2}}}        % ²
\DeclareUnicodeCharacter{00B3}{\ensuremath{^{3}}}        % ³
\DeclareUnicodeCharacter{00B5}{\ensuremath{\mu}}         % µ (micro)
\DeclareUnicodeCharacter{00B7}{\ensuremath{\cdot}}       % · (middle dot)
\DeclareUnicodeCharacter{00D7}{\ensuremath{\times}}      % ×
\DeclareUnicodeCharacter{00E4}{\"a}             % ä
\DeclareUnicodeCharacter{00E9}{\'e}             % é
\DeclareUnicodeCharacter{00F6}{\"o}             % ö
\DeclareUnicodeCharacter{00FC}{\"u}             % ü
\DeclareUnicodeCharacter{010C}{\v{C}}           % Č
\DeclareUnicodeCharacter{010D}{\v{c}}           % č
\DeclareUnicodeCharacter{0394}{\ensuremath{\Delta}}      % Δ
\DeclareUnicodeCharacter{03A3}{\ensuremath{\Sigma}}      % Σ
\DeclareUnicodeCharacter{03A9}{\ensuremath{\Omega}}      % Ω
\DeclareUnicodeCharacter{03B1}{\ensuremath{\alpha}}      % α
\DeclareUnicodeCharacter{03B2}{\ensuremath{\beta}}       % β
\DeclareUnicodeCharacter{03B3}{\ensuremath{\gamma}}      % γ
\DeclareUnicodeCharacter{03B4}{\ensuremath{\delta}}      % δ
\DeclareUnicodeCharacter{03B5}{\ensuremath{\varepsilon}} % ε
\DeclareUnicodeCharacter{03BA}{\ensuremath{\kappa}}      % κ
\DeclareUnicodeCharacter{03BB}{\ensuremath{\lambda}}     % λ
\DeclareUnicodeCharacter{03BC}{\ensuremath{\mu}}         % μ
\DeclareUnicodeCharacter{03BD}{\ensuremath{\nu}}         % ν
\DeclareUnicodeCharacter{03C0}{\ensuremath{\pi}}         % π
\DeclareUnicodeCharacter{03C1}{\ensuremath{\rho}}        % ρ
\DeclareUnicodeCharacter{03C3}{\ensuremath{\sigma}}      % σ
\DeclareUnicodeCharacter{03C4}{\ensuremath{\tau}}        % τ
\DeclareUnicodeCharacter{03C6}{\ensuremath{\varphi}}     % φ
\DeclareUnicodeCharacter{03C7}{\ensuremath{\chi}}        % χ
\DeclareUnicodeCharacter{03C8}{\ensuremath{\psi}}        % ψ
\DeclareUnicodeCharacter{03C9}{\ensuremath{\omega}}      % ω
\DeclareUnicodeCharacter{2013}{--}              % – (en dash)
\DeclareUnicodeCharacter{2014}{---}             % — (em dash)
\DeclareUnicodeCharacter{2018}{`}               % ‘
\DeclareUnicodeCharacter{2019}{'}               % ’
\DeclareUnicodeCharacter{201C}{``}              % “
\DeclareUnicodeCharacter{201D}{''}              % ”
\DeclareUnicodeCharacter{2070}{\ensuremath{^{0}}}        % ⁰
\DeclareUnicodeCharacter{2071}{\ensuremath{^{i}}}        % ⁱ
\DeclareUnicodeCharacter{2122}{\textsuperscript{TM}}     % ™
\DeclareUnicodeCharacter{2126}{\ensuremath{\Omega}}      % Ω (ohm sign)
\DeclareUnicodeCharacter{210F}{\ensuremath{\hbar}}       % ℏ
\DeclareUnicodeCharacter{2115}{\ensuremath{\mathbb{N}}}  % ℕ
\DeclareUnicodeCharacter{2102}{\ensuremath{\mathbb{C}}}  % ℂ
\DeclareUnicodeCharacter{211D}{\ensuremath{\mathbb{R}}}  % ℝ
\DeclareUnicodeCharacter{2124}{\ensuremath{\mathbb{Z}}}  % ℤ
\DeclareUnicodeCharacter{211A}{\ensuremath{\mathbb{Q}}}  % ℚ
\DeclareUnicodeCharacter{2192}{\ensuremath{\to}}         % →
\DeclareUnicodeCharacter{21D2}{\ensuremath{\Rightarrow}} % ⇒
\DeclareUnicodeCharacter{2200}{\ensuremath{\forall}}     % ∀
\DeclareUnicodeCharacter{2203}{\ensuremath{\exists}}     % ∃
\DeclareUnicodeCharacter{2208}{\ensuremath{\in}}         % ∈
\DeclareUnicodeCharacter{2209}{\ensuremath{\notin}}      % ∉
\DeclareUnicodeCharacter{2227}{\ensuremath{\wedge}}      % ∧
\DeclareUnicodeCharacter{2228}{\ensuremath{\vee}}        % ∨
\DeclareUnicodeCharacter{2229}{\ensuremath{\cap}}        % ∩
\DeclareUnicodeCharacter{222A}{\ensuremath{\cup}}        % ∪
\DeclareUnicodeCharacter{2245}{\ensuremath{\cong}}       % ≅
\DeclareUnicodeCharacter{2248}{\ensuremath{\approx}}     % ≈
\DeclareUnicodeCharacter{2260}{\ensuremath{\neq}}        % ≠
\DeclareUnicodeCharacter{2264}{\ensuremath{\leq}}        % ≤
\DeclareUnicodeCharacter{2265}{\ensuremath{\geq}}        % ≥
\DeclareUnicodeCharacter{22C5}{\ensuremath{\cdot}}       % ⋅
\DeclareUnicodeCharacter{2295}{\ensuremath{\oplus}}      % ⊕
\DeclareUnicodeCharacter{2297}{\ensuremath{\otimes}}     % ⊗

% ---- Mathematics & symbols ------------------------------------------
\usepackage{amsmath,amssymb,bm,mathrsfs,mathtools}
\usepackage{amsthm}             % theorem environments (load BEFORE hyperref)

% ---- Theorem environments (shared counter set, TECT canonical) -------
% All theorems / lemmas / propositions / corollaries / remarks share a
% single counter so cross-references read consistently across a paper.
% Definitions get a separate counter (definitions are numbered in their
% own sequence by editorial convention).
\newtheorem{theorem}{Theorem}
\newtheorem{lemma}[theorem]{Lemma}
\newtheorem{proposition}[theorem]{Proposition}
\newtheorem{corollary}[theorem]{Corollary}

\theoremstyle{remark}
\newtheorem{remark}[theorem]{Remark}
\newtheorem{example}[theorem]{Example}

\theoremstyle{definition}
\newtheorem{definition}[theorem]{Definition}
\newtheorem{conjecture}[theorem]{Conjecture}
\newtheorem{hypothesis}[theorem]{Hypothesis}

% ---- Tables / floats / graphics --------------------------------------
\usepackage{booktabs}           % \toprule / \midrule / \bottomrule
\usepackage{array}
\usepackage{graphicx}
\usepackage{enumitem}           % \begin{itemize}[leftmargin=...] options
\usepackage{hyphenat}           % \hyp{} for breakable hyphens in \texttt
\usepackage{seqsplit}           % \seqsplit{} for long unbreakable strings
\sloppy                         % allow stretchier inter-word spacing to
                                % reduce overfull \hbox warnings on long URLs
                                % and long \texttt tokens (paragraph-level).
\emergencystretch=3em           % last-resort hbox stretch budget

% ---- Cross-references (hyperref last among the standard packages,
%      cleveref AFTER hyperref) ----------------------------------------
\usepackage[%
  colorlinks=true,%
  linkcolor=blue,%
  citecolor=blue,%
  urlcolor=blue,%
  breaklinks=true,%
  bookmarks=true,%
  bookmarksnumbered=true%
]{hyperref}
\usepackage[capitalise,nameinlink]{cleveref}

% ---- TECT-canonical author block macros -----------------------------
% (defined here so every paper has the same author / affiliation style)
\newcommand{\tectauthor}{%
  \author{Jusang Lee}%
  \email{jtkor@outlook.com}%
  \affiliation{Independent researcher, Republic of Korea}%
}

% ---- TECT-canonical archive footer ----------------------------------
\newcommand{\tectarchivefooter}{%
  \par\medskip
  \noindent\small\textit{Canonical archive}:
  \url{https://github.com/TECT-OpenPhysics/TECT}.\par%
}

% ---- TECT TODO / AUDIT-FLAG / HONEST-SCOPE inline marks --------------
% Used in drafts to highlight pending audit items without disturbing
% the body text style. Suppress via \tectsuppressmarks (define empty).
\providecommand{\tectauditflag}[1]{\textcolor{red}{\textbf{[AUDIT-FLAG: #1]}}}
\providecommand{\tecttodo}[1]{\textcolor{orange}{\textbf{[TODO: #1]}}}
\providecommand{\tecthonestscope}[1]{\textit{Honest scope: #1.}}

% =====================================================================
% End of TECT canonical preamble
% =====================================================================

%   \begin{document}
%   ...
%   \end{document}
%
% Build (every paper): `pdflatex Paper-NN && pdflatex Paper-NN`
% (the second pass resolves \ref{} / \cref{} forward references).
% A bash/PowerShell wrapper that automates this is at
% Docs/papers/papers/_shared/build-paper.sh / build-paper.ps1.
% =====================================================================

% ---- Document class --------------------------------------------------
% PRD class option of revtex4-2 with [reprint,onecolumn] gives the same
% APS heading / by-line / abstract style as PRL but in a wider single
% column that handles long display equations and theorem environments
% without column-break artefacts. nofootinbib keeps citations out of
% footnotes. The 11pt option restores standard font metrics; with T1
% fontenc + lmodern this gives non-bitmap glyphs.
\documentclass[%
  aps,prd,reprint,onecolumn,nofootinbib,11pt,%
  amsmath,amssymb,floatfix,longbibliography%
]{revtex4-2}

% ---- Encoding & fonts ------------------------------------------------
\usepackage[T1]{fontenc}        % proper 8-bit font encoding (no font corruption)
\usepackage[utf8]{inputenc}     % allow UTF-8 source
\usepackage{lmodern}            % scalable Latin Modern (replaces bitmap CM)
\usepackage{textcomp}           % extra text symbols
\usepackage{microtype}          % subtle kerning improvements

% ---- UTF-8 → LaTeX character mappings --------------------------------
% Maps the Unicode characters that occur in TECT body text (legacy from
% the chat-content auto-archival rule, CLAUDE.md §4) to LaTeX-safe
% commands. Without these declarations, T1+inputenc rejects the chars
% with "Unicode character X not set up for use with LaTeX".
% Adding declarations centrally (here) is preferable to per-file edits.
\DeclareUnicodeCharacter{00A7}{\S}              % §  (section sign)
\DeclareUnicodeCharacter{00B2}{\ensuremath{^{2}}}        % ²
\DeclareUnicodeCharacter{00B3}{\ensuremath{^{3}}}        % ³
\DeclareUnicodeCharacter{00B5}{\ensuremath{\mu}}         % µ (micro)
\DeclareUnicodeCharacter{00B7}{\ensuremath{\cdot}}       % · (middle dot)
\DeclareUnicodeCharacter{00D7}{\ensuremath{\times}}      % ×
\DeclareUnicodeCharacter{00E4}{\"a}             % ä
\DeclareUnicodeCharacter{00E9}{\'e}             % é
\DeclareUnicodeCharacter{00F6}{\"o}             % ö
\DeclareUnicodeCharacter{00FC}{\"u}             % ü
\DeclareUnicodeCharacter{010C}{\v{C}}           % Č
\DeclareUnicodeCharacter{010D}{\v{c}}           % č
\DeclareUnicodeCharacter{0394}{\ensuremath{\Delta}}      % Δ
\DeclareUnicodeCharacter{03A3}{\ensuremath{\Sigma}}      % Σ
\DeclareUnicodeCharacter{03A9}{\ensuremath{\Omega}}      % Ω
\DeclareUnicodeCharacter{03B1}{\ensuremath{\alpha}}      % α
\DeclareUnicodeCharacter{03B2}{\ensuremath{\beta}}       % β
\DeclareUnicodeCharacter{03B3}{\ensuremath{\gamma}}      % γ
\DeclareUnicodeCharacter{03B4}{\ensuremath{\delta}}      % δ
\DeclareUnicodeCharacter{03B5}{\ensuremath{\varepsilon}} % ε
\DeclareUnicodeCharacter{03BA}{\ensuremath{\kappa}}      % κ
\DeclareUnicodeCharacter{03BB}{\ensuremath{\lambda}}     % λ
\DeclareUnicodeCharacter{03BC}{\ensuremath{\mu}}         % μ
\DeclareUnicodeCharacter{03BD}{\ensuremath{\nu}}         % ν
\DeclareUnicodeCharacter{03C0}{\ensuremath{\pi}}         % π
\DeclareUnicodeCharacter{03C1}{\ensuremath{\rho}}        % ρ
\DeclareUnicodeCharacter{03C3}{\ensuremath{\sigma}}      % σ
\DeclareUnicodeCharacter{03C4}{\ensuremath{\tau}}        % τ
\DeclareUnicodeCharacter{03C6}{\ensuremath{\varphi}}     % φ
\DeclareUnicodeCharacter{03C7}{\ensuremath{\chi}}        % χ
\DeclareUnicodeCharacter{03C8}{\ensuremath{\psi}}        % ψ
\DeclareUnicodeCharacter{03C9}{\ensuremath{\omega}}      % ω
\DeclareUnicodeCharacter{2013}{--}              % – (en dash)
\DeclareUnicodeCharacter{2014}{---}             % — (em dash)
\DeclareUnicodeCharacter{2018}{`}               % ‘
\DeclareUnicodeCharacter{2019}{'}               % ’
\DeclareUnicodeCharacter{201C}{``}              % “
\DeclareUnicodeCharacter{201D}{''}              % ”
\DeclareUnicodeCharacter{2070}{\ensuremath{^{0}}}        % ⁰
\DeclareUnicodeCharacter{2071}{\ensuremath{^{i}}}        % ⁱ
\DeclareUnicodeCharacter{2122}{\textsuperscript{TM}}     % ™
\DeclareUnicodeCharacter{2126}{\ensuremath{\Omega}}      % Ω (ohm sign)
\DeclareUnicodeCharacter{210F}{\ensuremath{\hbar}}       % ℏ
\DeclareUnicodeCharacter{2115}{\ensuremath{\mathbb{N}}}  % ℕ
\DeclareUnicodeCharacter{2102}{\ensuremath{\mathbb{C}}}  % ℂ
\DeclareUnicodeCharacter{211D}{\ensuremath{\mathbb{R}}}  % ℝ
\DeclareUnicodeCharacter{2124}{\ensuremath{\mathbb{Z}}}  % ℤ
\DeclareUnicodeCharacter{211A}{\ensuremath{\mathbb{Q}}}  % ℚ
\DeclareUnicodeCharacter{2192}{\ensuremath{\to}}         % →
\DeclareUnicodeCharacter{21D2}{\ensuremath{\Rightarrow}} % ⇒
\DeclareUnicodeCharacter{2200}{\ensuremath{\forall}}     % ∀
\DeclareUnicodeCharacter{2203}{\ensuremath{\exists}}     % ∃
\DeclareUnicodeCharacter{2208}{\ensuremath{\in}}         % ∈
\DeclareUnicodeCharacter{2209}{\ensuremath{\notin}}      % ∉
\DeclareUnicodeCharacter{2227}{\ensuremath{\wedge}}      % ∧
\DeclareUnicodeCharacter{2228}{\ensuremath{\vee}}        % ∨
\DeclareUnicodeCharacter{2229}{\ensuremath{\cap}}        % ∩
\DeclareUnicodeCharacter{222A}{\ensuremath{\cup}}        % ∪
\DeclareUnicodeCharacter{2245}{\ensuremath{\cong}}       % ≅
\DeclareUnicodeCharacter{2248}{\ensuremath{\approx}}     % ≈
\DeclareUnicodeCharacter{2260}{\ensuremath{\neq}}        % ≠
\DeclareUnicodeCharacter{2264}{\ensuremath{\leq}}        % ≤
\DeclareUnicodeCharacter{2265}{\ensuremath{\geq}}        % ≥
\DeclareUnicodeCharacter{22C5}{\ensuremath{\cdot}}       % ⋅
\DeclareUnicodeCharacter{2295}{\ensuremath{\oplus}}      % ⊕
\DeclareUnicodeCharacter{2297}{\ensuremath{\otimes}}     % ⊗

% ---- Mathematics & symbols ------------------------------------------
\usepackage{amsmath,amssymb,bm,mathrsfs,mathtools}
\usepackage{amsthm}             % theorem environments (load BEFORE hyperref)

% ---- Theorem environments (shared counter set, TECT canonical) -------
% All theorems / lemmas / propositions / corollaries / remarks share a
% single counter so cross-references read consistently across a paper.
% Definitions get a separate counter (definitions are numbered in their
% own sequence by editorial convention).
\newtheorem{theorem}{Theorem}
\newtheorem{lemma}[theorem]{Lemma}
\newtheorem{proposition}[theorem]{Proposition}
\newtheorem{corollary}[theorem]{Corollary}

\theoremstyle{remark}
\newtheorem{remark}[theorem]{Remark}
\newtheorem{example}[theorem]{Example}

\theoremstyle{definition}
\newtheorem{definition}[theorem]{Definition}
\newtheorem{conjecture}[theorem]{Conjecture}
\newtheorem{hypothesis}[theorem]{Hypothesis}

% ---- Tables / floats / graphics --------------------------------------
\usepackage{booktabs}           % \toprule / \midrule / \bottomrule
\usepackage{array}
\usepackage{graphicx}
\usepackage{enumitem}           % \begin{itemize}[leftmargin=...] options
\usepackage{hyphenat}           % \hyp{} for breakable hyphens in \texttt
\usepackage{seqsplit}           % \seqsplit{} for long unbreakable strings
\sloppy                         % allow stretchier inter-word spacing to
                                % reduce overfull \hbox warnings on long URLs
                                % and long \texttt tokens (paragraph-level).
\emergencystretch=3em           % last-resort hbox stretch budget

% ---- Cross-references (hyperref last among the standard packages,
%      cleveref AFTER hyperref) ----------------------------------------
\usepackage[%
  colorlinks=true,%
  linkcolor=blue,%
  citecolor=blue,%
  urlcolor=blue,%
  breaklinks=true,%
  bookmarks=true,%
  bookmarksnumbered=true%
]{hyperref}
\usepackage[capitalise,nameinlink]{cleveref}

% ---- TECT-canonical author block macros -----------------------------
% (defined here so every paper has the same author / affiliation style)
\newcommand{\tectauthor}{%
  \author{Jusang Lee}%
  \email{jtkor@outlook.com}%
  \affiliation{Independent researcher, Republic of Korea}%
}

% ---- TECT-canonical archive footer ----------------------------------
\newcommand{\tectarchivefooter}{%
  \par\medskip
  \noindent\small\textit{Canonical archive}:
  \url{https://github.com/TECT-OpenPhysics/TECT}.\par%
}

% ---- TECT TODO / AUDIT-FLAG / HONEST-SCOPE inline marks --------------
% Used in drafts to highlight pending audit items without disturbing
% the body text style. Suppress via \tectsuppressmarks (define empty).
\providecommand{\tectauditflag}[1]{\textcolor{red}{\textbf{[AUDIT-FLAG: #1]}}}
\providecommand{\tecttodo}[1]{\textcolor{orange}{\textbf{[TODO: #1]}}}
\providecommand{\tecthonestscope}[1]{\textit{Honest scope: #1.}}

% =====================================================================
% End of TECT canonical preamble
% =====================================================================


\begin{document}

\title{A four-sector cancellation programme note for the cosmological
constant in TECT \\
\normalsize{\textit{(NOT a Pillar-11 closure paper; the present
draft fails its own pre-registered falsification gate as written ---
internal numerical / gate inconsistency flagged for revision.)}}}

\author{Jusang Lee}
\email{jtkor@outlook.com}
\affiliation{Independent researcher, Republic of Korea}

\date{\today}

\begin{abstract}
We outline a candidate four-sector cancellation programme for the
cosmological constant within the Topological Energy Condensate Theory
(TECT). The BCC topological condensate is decomposed into four
candidate independent sectors: (i) the BCC order-parameter
self-interaction energy, (ii) the monopole sector (captured by the
topological Chern class $c_2$), (iii) the Dirac-fermion zero-mode
sector, and (iv) the vortex sector (defect core energy).

\textbf{Internal-inconsistency flag (AUDIT-2026-05-02-Wave7,
Math314-AddD)}: as drafted, this paper claims a near-cancellation
residual $|\sum_i E_i^{\rm vac}| < 10^{-10}\hbar\omega_0$ and
identifies $\Lambda_{\rm TECT} \approx 10^{-120}\,M_{\rm Planck}^4$
in §Mechanism, but the §Numerical-result computation gives
$|\Lambda_{\rm TECT}|/M_{\rm Planck}^4 \sim 10^{-118}\text{ to }10^{-119}$,
which DOES NOT PASS the pre-registered falsification gate
$|\Lambda_{\rm TECT}|/M_{\rm Planck}^4 < 10^{-120}$. The paper
internally fails its own falsification gate. The Pillar-11 closure
claim and the ``T6 PROVED CONDITIONAL'' tier statement are
therefore RETRACTED at the present canonical tier; the document is
recast as a candidate programme note pending (a) repair of the
numerical-vs-gate inconsistency, (b) honest re-statement of the
expected cancellation tier, and (c) verdict-shell scheduling per the
canonical Math3xx framework. The current canonical tier of
Pillar~11 is recorded in \texttt{Docs/status/TOE-FACT-SHEET.md}.
\end{abstract}

\maketitle

% =====================================================================
\section{Introduction}\label{sec:intro}
% =====================================================================
The cosmological constant problem is one of the most acute tensions
in theoretical physics: the observed value
$\Lambda_{\mathrm{obs}} \approx 10^{-120}\,M_{\mathrm{Planck}}^4$
(type-Ia supernovae, CMB, BAO) is many orders of magnitude smaller
than naive quantum-field-theory vacuum-energy estimates.

TECT \emph{proposes} a candidate four-sector cancellation programme
in which the cosmological constant arises as the residual of
near-cancellation among four topologically-distinguished sectors of
the BCC condensate vacuum energy. At the present canonical tier the
programme \emph{conjectures} that the cancellation is enforced by
lattice topological selection rules rather than by a fine-tuning of
couplings; the conjecture is \emph{not} discharged at theorem level
in this paper.

The present paper \emph{records} the four-sector decomposition, the
candidate cancellation condition, the explicit numerical operating
point, and the \emph{pre-registered falsification gate that the
present draft fails as written} (residual is observed at
$10^{-118}$ to $10^{-119}$ in $M_{\mathrm{Planck}}^4$ units, while
the gate requires $< 10^{-120}$). The Pillar~11 closure claim is
\emph{not} made here; the document is a programme note pending (i)
repair of the numerical-vs-gate inconsistency, (ii) honest
re-statement of the expected cancellation tier, and (iii)
verdict-shell scheduling per the canonical Math3xx framework.

% =====================================================================
\section{The four-sector decomposition}
% =====================================================================

The TECT vacuum energy (at one-loop order) decomposes into four sectors:

\subsection{Sector 1: BCC order-parameter self-interaction}

The Brazovskii free energy yields a one-loop effective potential:
\begin{equation}
E_1^{\rm vac} = \frac{1}{2} \int \frac{d^3\mathbf{k}}{(2\pi)^3} \omega_\mathbf{k},
\end{equation}
where $\omega_\mathbf{k}$ is the one-particle dispersion of the BCC order parameter.
In the broken phase, $\omega_\mathbf{k} = \sqrt{m^{*2} + (k/m^*)^2}$ (with $m^*$
the mass gap from Pillar 1). The integral is regularized via dimensional
regularization and renormalized mass (Math58-v3 framework):
\begin{equation}
E_1^{\rm vac} \sim \lambda^{1/3} M_X^4 + \text{(log terms)}.
\end{equation}

\subsection{Sector 2: Monopole sector (topological Chern class)}

The BCC lattice admits topological monopole excitations, whose vacuum contribution
arises from the loop integral over monopole propagation. The Chern-class index
$c_2(E)$ of the monopole gauge bundle constrains the spectrum. Under the gauge
choice $c_1 = 0$ (Math191), the Chern number becomes:
\begin{equation}
c_2(E) = \int_X c_2 = -40 \quad \text{(on } \Sigma_0 \text{)}.
\end{equation}
The vacuum energy from this sector is:
\begin{equation}
E_2^{\rm vac} = -\gamma^{2/3} M_X^4 + \text{(log terms)},
\end{equation}
with the opposite sign from Sector 1.

\subsection{Sector 3: Dirac-fermion zero-mode sector}

The protected chiral zero mode (Pillar 5, via index theorem) contributes a
fermionic zero-point energy. Since it is a single zero mode per BCC unit cell,
the total fermionic contribution is:
\begin{equation}
E_3^{\rm vac} = \frac{1}{2} \hbar v_F q_0 \quad \text{(one Dirac zero mode per cell)}.
\end{equation}
In natural units and volume-normalized, this becomes:
\begin{equation}
E_3^{\rm vac} = \alpha_D M_X^4 + \text{(log terms)},
\end{equation}
with $\alpha_D$ a dimensionless coefficient depending on the Fermi velocity.

\subsection{Sector 4: Vortex sector (defect core energy)}

Topological defects (vortices) in the BCC order parameter carry core energy.
The vacuum loop integral over vortex pairs yields:
\begin{equation}
E_4^{\rm vac} = -(\alpha_1 \lambda^{1/3} + \alpha_2 \gamma^{2/3}) M_X^4 + \text{(log terms)},
\end{equation}
with $\alpha_1, \alpha_2$ dimensionless coefficients.

% =====================================================================
\section{The candidate cancellation condition (programme statement,
not theorem)}\label{sec:cancellation}
% =====================================================================

The four-sector sum reads
\begin{equation}\label{eq:Etotal}
  E_{\mathrm{vac}}^{\mathrm{total}}
  \;=\; \big(E_1
    - \alpha_1 \lambda^{1/3}
    + \alpha_D
    - \gamma^{2/3}\,E_2\big)\,M_X^4
  + \text{(log)}.
\end{equation}
The candidate near-cancellation condition, at the tree level before
renormalisation, is
\begin{equation}\label{eq:cancel-cond}
  1 \;-\; \alpha_1 \lambda^{1/3} \;+\; \alpha_D \;-\; \gamma^{2/3}
  \;\approx\; 0.
\end{equation}

\begin{remark}[The cancellation \eqref{eq:cancel-cond} is a candidate
programme statement, NOT a theorem at the present canonical
tier]\label{rmk:cancel-not-theorem}
The earlier draft of this paper asserted that
\eqref{eq:cancel-cond} is an ``automatic consequence'' of the
topological selection rules and that Math58-v7 ``proves'' the
coefficients $\{\alpha_1,\alpha_2,\alpha_D\}$ are forced by the
lattice geometry and Dirac-index invariants. That assertion is
\emph{programme}, not \emph{theorem}. Even within Math58-v7, the
coefficient identification is performed under the load-bearing
hypotheses (H1)--(H3) listed in \S\ref{sec:hypotheses-H}, and the
matching of the resulting near-cancellation residual to
$\Lambda_{\mathrm{obs}}$ requires the verdict-shell scheduling of
the canonical Math3xx framework. At the present canonical tier
\eqref{eq:cancel-cond} is a \textbf{T2 \textsc{conjecture}}, not a
\textbf{T7 \textsc{proved}} theorem; the corresponding Pillar~11
status is \textbf{T4 \textsc{strong evidence}}.
\end{remark}

If \eqref{eq:cancel-cond} held exactly, the residual vacuum energy
would arise from one-loop corrections to the cancellation:
\begin{equation}\label{eq:Eres}
  E_{\mathrm{vac}}^{\mathrm{residual}}
  \;\sim\; \hbar\,\frac{\lambda^2}{16\pi^2}\,M_X^4
  + \text{(higher-order)}.
\end{equation}
At the Math58-v7 operating point with $\lambda \sim 10^{-3}$ at the
renormalisation scale, \eqref{eq:Eres} estimates
\begin{equation}\label{eq:Eres-num}
  E_{\mathrm{vac}}^{\mathrm{residual}}
  \;\sim\; 10^{-10}\,\hbar\,M_X^4.
\end{equation}
The estimate \eqref{eq:Eres-num} is a programme-level scaling
estimate; it is \emph{not} a constant-bound theorem.

% =====================================================================
\section{Candidate connection to the observed cosmological constant
(target value, NOT verified prediction)}\label{sec:hypotheses-H}
% =====================================================================

The TECT vacuum energy would be related to the cosmological constant
via the Einstein equation,
\begin{equation}\label{eq:einstein-Lambda}
  \Lambda_{\mathrm{TECT}}
  \;=\; \frac{8\pi G}{3}\,\rho_{\mathrm{vac}}
  \;=\; \frac{8\pi G}{3}\,
        \frac{E_{\mathrm{vac}}^{\mathrm{residual}}}{V}.
\end{equation}
Substituting \eqref{eq:Eres-num} and the TECT-derived constants
$G, c, a_{\mathrm{BCC}}$ from the prior pillars yields the
\emph{target value}
\begin{equation}\label{eq:Lambda-target}
  \Lambda_{\mathrm{TECT}}^{\mathrm{target}}
  \;\approx\; 10^{-120}\,M_{\mathrm{Planck}}^4,
\end{equation}
which would match the observed value
($\Omega_\Lambda \approx 0.73$). The actual numerical evaluation of
the four-sector residual at the Math58-v7 operating point gives
$|\Lambda_{\mathrm{TECT}}|/M_{\mathrm{Planck}}^4 \sim 10^{-118}$ to
$10^{-119}$ (\S\ref{sec:numerics}), which differs from
\eqref{eq:Lambda-target} by $1$--$2$ orders of magnitude. The
target--evaluation gap is the canonical falsification-gate failure
recorded in Remark~\ref{rmk:gate-fail}.

The candidate programme rests on three explicit load-bearing
hypotheses.
\begin{enumerate}
  \item \textbf{(H1) Chern-class exactness (Math191).} The gauge
    choice $c_1 = 0$ leaves the Chern number $c_2 = -40$ invariant
    and the cancellation \emph{schematic} as written. Conditional on
    the mathematical exactness of the Chern-class formula on the
    BCC lattice.
  \item \textbf{(H2) Vortex topology stability (Math58-v4/v5).}
    Topological defects on the BCC lattice preserve their winding
    numbers under small perturbations. Topological property; binary.
  \item \textbf{(H3) Lattice--continuum matching (Math58-v6/v7).}
    The continuum limit of the four-sector loop integrals matches
    the discrete lattice calculation to within $\sim 10^{-12}$
    relative error. Conditional on the discretisation scheme;
    Math58-v8 tests the continuum limit.
\end{enumerate}

% =====================================================================
\section{Numerical evaluation and the
falsification-gate failure as written}\label{sec:numerics}
% =====================================================================

Math58-v7 §Q5 executes a numerical evaluation of the four-sector
cancellation \eqref{eq:cancel-cond}. The operating point is
\begin{equation}\label{eq:opt-pt}
  (\lambda,\gamma,Y,q_0,a_{\mathrm{BCC}})
  \;=\; (-0.0394,-0.0188,280.2,14.32,4\sqrt{\pi}\,\ell_P).
\end{equation}
The evaluated residual is
\begin{equation}\label{eq:Lambda-num}
  \frac{|\Lambda_{\mathrm{TECT}}|}{M_{\mathrm{Planck}}^4}
  \;\approx\; 10^{-118}\;\text{ to }\;10^{-119}.
\end{equation}

\paragraph{Pre-registered falsification gate.}
\begin{equation}\label{eq:fgate}
  \frac{|\Lambda_{\mathrm{TECT}}|}{M_{\mathrm{Planck}}^4}
  \;<\; 10^{-120}.
\end{equation}

\begin{remark}[The present draft FAILS the falsification gate
\eqref{eq:fgate} as written]\label{rmk:gate-fail}
\eqref{eq:Lambda-num} is in the range $10^{-118}$ to $10^{-119}$,
which exceeds the gate threshold $10^{-120}$ of \eqref{eq:fgate} by
$1$--$2$ orders of magnitude. The pre-registered falsification gate
is therefore \emph{not passed} by the present draft. This is a
self-failure recorded honestly per CLAUDE.md~§6.3.4 quantitative
sanity check. Pillar~11 \emph{cannot} be promoted to T6 on the
basis of this paper; the canonical tier reverts to (or remains at)
\textbf{T4 \textsc{strong evidence}} with this gate-failure
recorded as part of the verdict-shell scheduling.
\end{remark}

The 2026-06-01 verdict scheduled in
\texttt{PHASE\_8\_TO\_14\_PLAN.md} §11 must therefore consume the
gate-failure as a condition of promotion: either (a) the operating
point \eqref{eq:opt-pt} or the cancellation coefficient identification
of \S\ref{sec:cancellation} must be revised to drive the residual
below $10^{-120}$, or (b) the gate \eqref{eq:fgate} itself must be
re-registered at a different threshold reflecting the canonical
near-cancellation expectation. Until either option is exercised and
audited, Pillar~11 status remains \textbf{T4 \textsc{strong
evidence}}.

% =====================================================================
\section{Conclusion}\label{sec:conclusion}
% =====================================================================

The TECT four-sector cancellation programme of \S\ref{sec:cancellation}
is a candidate mechanism for the smallness of the cosmological
constant within the BCC condensate. It is recorded here as a
\textbf{T2 \textsc{conjecture}} on the cancellation, supported by an
order-of-magnitude estimate \eqref{eq:Eres-num} of the residual. The
target value \eqref{eq:Lambda-target} matches the observed value
$\Lambda_{\mathrm{obs}}$ in order of magnitude, but the actual
numerical evaluation \eqref{eq:Lambda-num} differs from the
pre-registered falsification gate \eqref{eq:fgate} by one to two
orders of magnitude. The Pillar~11 canonical tier remains
\textbf{T4 \textsc{strong evidence}}; promotion to T6 \textsc{proved
conditional} is contingent on (a) the gate-failure repair tracked in
Remark~\ref{rmk:gate-fail}, and (b) the verdict-shell scheduling per
the canonical Math3xx framework. This document is not a Pillar~11
closure paper; it is a programme note in which the pre-registered
gate is honestly reported as failed.

\begin{thebibliography}{99}

\bibitem{TECTRepo}
J. Lee, \textit{Topological Energy Condensate Theory (TECT) Repository},
\texttt{https://github.com/...} (2026).

\bibitem{Math58v7}
J. Lee, ``Pillar 11 Dirac-sector tightening,''
\textit{Math58-v7}, 2026-04-25.

\bibitem{Math58v7AddC}
J. Lee, ``Pillar 11 promotion to PROVED CONDITIONAL,''
\textit{Math58-v7-AddC}, 2026-04-25.

\bibitem{Math191}
J. Lee, ``Gauge choice $c_1 = 0$ canonical form,''
\textit{Math191}, 2026-04-25.

\bibitem{Math58v7AddA}
J. Lee, ``PV scheme adversarial audit,''
\textit{Math58-v7-Addendum-A}, 2026-04-25.

\bibitem{Math58v8}
J. Lee, ``Pillar 11 continuum-limit closure,''
\textit{Math58-v8}, 2026-04-26.

\end{thebibliography}

\end{document}
